\begin{glyph}{Dry (干)}{dry}\label{glyph:dry}
Dry.  Ebb.
\end{glyph}

\begin{comps}
\comprtwocone & 干 \\\hline
\comprthreecone & Dry/Shield \\\hline
\end{comps}

\begin{remglyph}
Part of the 漢字 radicals (\#{51}).\\\hline
\end{remglyph}

\begin{prons}
\pronsrtwocone & --- & \textbf{ひ、ほ (hi, ho)}\\\hline
\pronsrthreecone & カン (KAN) & --- \\\hline
\end{prons}

\begin{remprons}
    \hooglig{Dry} \comkun{heaps} of \hooglig{dried} fruit garnered by
    \ucomon{cunning} \comkun{hobbyists}.\\\hline
\end{remprons}

\begin{somme}
    干る (intr.) & \textit{to dry (on its own)} & ひ{\middel}る \mbox{(hi{\middel}ru)}\\\hline
    干す (tr.) & \textit{to dry (something)} & ほ{\middel}す \mbox{(ho{\middel}su)}\\\hline
    {\diff}日干し & sun + dry = \textit{sun-dried} & ひ　ぼ{\middel}し \mbox{(hi bo{\middel}shi)}\\\hline
    干上がる & dry + upper = \textit{to dry up} & ひ　あ{\middel}がる \mbox{(hi a{\middel}garu)}\\\hline
    \notyet{干物} & dry + thing = \textit{dried fish} & ひ{\middel}もの \mbox{(hi{\middel}mono)}\\\hline
\end{somme}

\begin{sentence}
\ruby{\underline{干物}}{ひもの}\ruby{を}{お}\ruby{\underline{買}}{か}\underline{って} \ruby{\underline{下}}{くだ}\underline{さい}。\\\hline
hi{\middel}mono o kat{\middel}te kuda{\middel}sai.\\\hline
(dried fish) (buy) (please)\\\hline
=\textit{Please buy dried fish.}\\\hline
\end{sentence}
\end{multicols}
\vhrulefill{2pt}
%%--------------------------------------------------------------------
%%--------------------------------------------------------------------
\begin{multicols}{2}
\begin{glyph}{Gold (金)}{gold}\label{glyph:gold}
Depending on the context in which it appears, this common character signifies either ``gold'', ``metal'', or ``money''.\\\hline
It clearly means ``gold'', for instance, when referring to the status of an Olympic champion.\\\hline
When functioning as a component, however, it hints that the 漢字 involved has some connection to metal in general.
\end{glyph}

\begin{comps}
\comprtwocone & 金\\\hline
\comprthreecone &　Gold\\\hline
\end{comps}

\begin{remglyph}
Part of the Kanji radicals (\#{167}).\\\hline
\end{remglyph}

\begin{prons}
\pronsrtwocone & \textbf{キン (KIN)} & \textbf{かね (kane)}\\\hline
\pronsrthreecone & コン (KON) & かな (kana) \\\hline
\end{prons}

\begin{remprons}
The \hooglig{golden} \ucomkun{canary} \comkun{cane} were stolen by \comon{kindred} \ucomon{con} artists.\\\hline
\end{remprons}

\begin{somme}
    金 & \textit{gold; metal; money} & キン \mbox{(KIN)}\\\hline
    金 & \textit{metal; money} & かね \mbox{(kane)}\\\hline
    金山 & gold + mountain = \textit{gold mine} & キン{\middel}ザン \mbox{(KIN{\middel}ZAN)}\\\hline
    \notyet{年金} & year + gold = \textit{annuity/pension} & ネン{\middel}キン \mbox{(NEN{\middel}KIN)}\\\hline
    \notyet{金曜日} & gold + day of the week + sun (day) = \textit{Friday} & キン{\middel}ヨウ{\middel}び \mbox{(KIN{\middel}Y\={O}{\middel}bi)}\\\hline
    \notyet{金星} & gold + star = \textit{Venus (planet)} & キン{\middel}セイ \mbox{(KIN{\middel}SEI)}\\\hline
\end{somme}

\begin{sentence}
\ruby{\underline{内田}}{うちだ}\underline{さん}\ruby{は}{わ}\underline{あまり} \underline{お}\ruby{\underline{金}}{かね}が\underline{ない}。\\\hline
Uchi{\middel}da-san wa amari o{\middel}kane ga nai.\\\hline
(Uchida-san) (not much) (money) (not)\\\hline
=\textit{Uchida-san doesn't have much money.}\\\hline
\end{sentence}
\end{multicols}
%%--------------------------------------------------------------------
%%--------------------------------------------------------------------
\begin{multicols}{2}
\begin{glyph}{Words (語)}{words}\label{glyph:words}
Words.\\\hline
This 漢字 will become familiar as the character is used to indicate languages.
\end{glyph}

\begin{comps}
\comprtwocone & 言 + 五 + 口 \\\hline
\comprthreecone & Say + Five + Mouth \\\hline
\end{comps}

\begin{remglyph}
\hooglig{Wording} the \remcom{five} \remcom{vampire} \remcom{sayings}.\\\hline
\end{remglyph}

\begin{prons}
\pronsrtwocone & \textbf{ゴ (GO)} & \textbf{かた (kata)}\\\hline
\pronsrthreecone & --- & かたり (katari) \\\hline
\end{prons}

\begin{remprons}
\hooglig{Word} \comkun{catalogues} \comon{gobbled} down to \ucomkun{cut a
    ribbon} at the finishline.\\\hline
\end{remprons}

\begin{somme}
\rye{3}{Note the importance of long and short vowel sounds in Japanese by looking closely at the second and third examples below.}
    語る & \textit{to talk} & かた{\middel}る \mbox{(kata{\middel}ru)}\\\hline
    口語 & mouth words = \textit{colloquial language} & コウ{\middel}ゴ \mbox{(K\={O}{\middel}GO)}\\\hline
    古語 & old + words = \textit{archaic words} & コ{\middel}ゴ \mbox{(KO{\middel}GO)}\\\hline
    日本語 & sun + main + words = \textit{Japanese (language)} & ニ{\middel}ホン{\middel}ゴ \mbox{(NI{\middel}HON{\middel}GO)}\\\hline
    言語 & say + words = \textit{language} & ゲン{\middel}ゴ \mbox{(GEN{\middel}GO)}\\\hline
    \notyet{英語} & English + words = \textit{English (language)} & エイ{\middel}ゴ \mbox{(EI{\middel}GO)}\\\hline
    \notyet{語学} & words + study = \textit{linguistics} & ゴ{\middel}ガク \mbox{(GO{\middel}GAKU)}\\\hline
\end{somme}

\begin{sentence}
\underline{あなた}の\ruby{\underline{日本語}}{ニホンゴ}\ruby{は}{わ}\ruby{\underline{上手}}{ジョウず}ですよ。\\\hline
anata no NI{\middel}HON{\middel}GO wa J\={O}{\middel}zu desu yo.\\\hline
(you) (Japanese) (good at)\\\hline
=\textit{Your Japanese is good.}\\\hline
\end{sentence}
\end{multicols}
\vhrulefill{2pt}
%--------------------------------------------------------------------
%--------------------------------------------------------------------
\begin{multicols}{2}
\begin{glyph}{Gentleman (士)}{gentleman}\label{glyph:gentleman}
Gentleman.\\\hline
A man who has attained status in a military or academic field.
\end{glyph}

\begin{comps}
\comprtwocone & 士 \\\hline
\comprthreecone & Gentleman \\\hline
\end{comps}

\begin{remglyph}
Part of the 漢字 radicals (\#{33}).\\\hline
\end{remglyph}

\begin{prons}
\pronsrtwocone & \textbf{シ (SHI)} & ---\\\hline
\pronsrthreecone & --- & --- \\\hline
\end{prons}

\begin{remprons}
\hooglig{Gentlemen} \comon{shield} other people.\\\hline
\end{remprons}

%{\middel}
\begin{somme}
    力士 & strength + gentleman = \textit{sumo wrestler} & リキ{\middel}シ \mbox{(RIKI{\middel}SHI)}\\\hline
    \notyet{士気} & gentleman + spirit = \textit{morale} & シ{\middel}キ \mbox{(SHI{\middel}KI)}\\\hline
    \notyet{工学士} & craft + study + gentleman = \textit{Bachelor of Engineering} & コウ{\middel}ガク{\middel}シ \mbox{(K\={O}{\middel}GAKU{\middel}SHI)}\\\hline
\end{somme}

\begin{sentence}
\underline{あの} \ruby{\underline{力士}}{リキシ}\ruby{は}{わ}\underline{とても} \ruby{\underline{大}}{おお}\underline{きい} \underline{です}よ。\\\hline
ano RIKI{\middel}SHI wa totemo \={o}{\middel}kii desu yo.\\\hline
(that) (sumo wrestler) (very) (large) (is)\\\hline
=\textit{That sumo wrestler is really large.}\\\hline
\end{sentence}
\end{multicols}
\vspace*{\fill}
%--------------------------------------------------------------------
%--------------------------------------------------------------------
\pagebreak
\vspace*{\fill}
\begin{multicols}{2}
\begin{glyph}{Morning (朝)}{morning}\label{glyph:morning}
Morning.\\\hline
A secondary meaning relates to ``royal dynasties''.
\end{glyph}

\begin{comps}
\comprtwocone & 十 + 早 + 月 \\\hline
\comprthreecone & Ten + Early (\ref{glyph:early}) + Moon\\\hline
\end{comps}

\begin{remglyph}
While the \hooglig{morning} is still \remcom{early}, the \remcom{moon} winks at
    the \remcom{scarecrow}.\\\hline
\end{remglyph}

\begin{prons}
\pronsrtwocone & \textbf{チョウ (CH\={O})} & \textbf{あさ (asa)}\\\hline
\pronsrthreecone & --- & --- \\\hline
\end{prons}

\begin{remprons}
\hooglig{Morning} \comon{chores} must be finished \comkun{asap}.\\\hline
\end{remprons}

%{\middel}
% &  ()\\\hline
\begin{somme}
    朝 & \textit{morning} & あさ \mbox{(asa)}\\\hline
    朝日 & morning + sun = \textit{morning sun} & あさ{\middel}ひ \mbox{(asa{\middel}hi)}\\\hline
    王朝 & king + morning = \textit{dynasty} & オウ{\middel}チョウ \mbox{(\={O}{\middel}CH\={O})}\\\hline
    早朝 & early + morning = \textit{early morning} & ソウ{\middel}チョウ \mbox{(S\={O}{\middel}CH\={O})}\\\hline
    \notyet{毎朝} & every + morning = \textit{every morning} & マイ{\middel}あさ \mbox{(MAI{\middel}asa)}\\\hline
\end{somme}

\begin{irrr}
    \notyet{今朝} & now + morning = \textit{this morning} & けさ (kesa)\\\hline
\end{irrr}

\begin{sentence}
\ruby{\underline{上田}}{うえだ}\underline{さん}\ruby{は}{わ}\ruby{\underline{朝}}{あさ}が\ruby{\underline{大好}}{ダイす}\underline{き}です。\\\hline
Ue{\middel}da-san wa asa ga DAI{\middel}su{\middel}ki desu.\\\hline
(Ueda-san) (morning) (really likes)\\\hline
=\textit{Ueda-san really likes mornings.}\\\hline
\end{sentence}
\end{multicols}
\vhrulefill{2pt}
%--------------------------------------------------------------------
%--------------------------------------------------------------------
\begin{multicols}{2}
\begin{glyph}{Blue (青)}{blue}\label{glyph:blue}
Blue.\\\hline
This character also applies to a few objects that Westerners invariably consider green.  E.g. \textit{blue traffic light}, \textit{blue vegetables}, amongst other things.\\\hline
The sense of ``youth'' and ``immaturity'' is tied up with the idea of ``green''.
\end{glyph}

\begin{comps}
\comprtwocone & 青 \\\hline
\comprthreecone & Blue \\\hline
\end{comps}

\begin{remglyph}
Part of the 漢字 radicals (\#{174}).\\\hline
\end{remglyph}

\begin{prons}
\pronsrtwocone & \textbf{セイ (SEI)} & \textbf{あお (ao)}\\\hline
\pronsrthreecone & ショウ (SH\={O}) &  \\\hline
\end{prons}

\begin{remprons}
\hooglig{Blue} \ucomon{shore} \comon{sailing} really pumps the \comkun{aorta}.\\\hline
\end{remprons}

%{\middel}
% &  ()\\\hline
\begin{somme}
    青 & \textit{blue; green (noun)} & あお \mbox{(ao)}\\\hline
    青い & \textit{blue; green (adjective)} & あお{\middel}い \mbox{(ao{\middel}i)}\\\hline
    青春 & blue + spring = \textit{youth} & セイ{\middel}シュン \mbox{(SEI{\middel}SHUN)}\\\hline
    {\weird}\notyet{青物} & blue + thing = \textit{green vegetables} & あお{\middel}もの \mbox{(ao{\middel}mono)}\\\hline
    \notyet{青空} & blue + empty = \textit{blue sky} & あお{\middel}ぞら \mbox{(ao{\middel}zora)}\\\hline
\end{somme}

\begin{sentence}
\ruby{\underline{山田}}{やまだ}\underline{さん}の\ruby{\underline{家}}{いえ}\ruby{は}{わ}\ruby{\underline{青}}{あお}\underline{い}です。\\\hline
Yama{\middel}da-san no ie wa ao{\middel}i desu.\\\hline
(Yamada-san) (house) (blue)\\\hline
=\textit{Yamada-san's house is blue.}\\\hline
\end{sentence}
\end{multicols}
\vspace*{\fill}
%%--------------------------------------------------------------------
%%--------------------------------------------------------------------
\pagebreak
\vspace*{\fill}
\begin{multicols}{2}
\begin{glyph}{Earth (土)}{earth}\label{glyph:earth}
Earth as in ``dirt'', not the planet.
\end{glyph}

\begin{comps}
\comprtwocone & 土 \\\hline
\comprthreecone & Earth \\\hline
\end{comps}

\begin{remglyph}
Part of the 漢字 radicals (\#{32}).\\\hline
\end{remglyph}

\begin{prons}
\pronsrtwocone & \textbf{ド (DO)} & \textbf{つち (tsuchi)}\\\hline
\pronsrthreecone & ト (TO) & --- \\\hline
\end{prons}

\begin{remprons}
\hooglig{Earth} \comon{dots} made by the \ucomon{top} \comkun{tsuchikage}.\\\hline
{\warn}Tsuchikage is a fictional アニメ character in ナルト.\\\hline
\end{remprons}

%{\middel}
% &  ()\\\hline
\begin{somme}
    土 & \textit{earth} & つち \mbox{(tsuchi)}\\\hline
    土手 & earth + hand = \textit{embankment} & ド{\middel}て \mbox{(DO{\middel}te)}\\\hline
    {\diff}本土 & main + earth = \textit{mainland} & ホン{\middel}ド \mbox{(HON{\middel}DO)}\\\hline
    土木 & earth + tree = \textit{Civil Engineering} & ド{\middel}ボク \mbox{(DO{\middel}BOKU)}\\\hline
    \notyet{土星} & earth + star = \textit{Saturn} & ド{\middel}セイ \mbox{(DO{\middel}SEI)}\\\hline
    \notyet{土曜日} & earth + day of the week + sun (day) = \textit{Saturday} & ド{\middel}ヨウ{\middel}び \mbox{(DO{\middel}Y\={O}{\middel}bi)}\\\hline
\end{somme}

\begin{irrr}
    \notyet{土産} & earth + produce = \textit{souvenir} & みやげ (miyage)\\\hline
\end{irrr}

\begin{sentence}
\ruby{\underline{川}}{かわ}の\ruby{\underline{土手}}{ドて}\ruby{を}{お}\ruby{\underline{上}}{あ}\underline{がる}と\ruby{\underline{水田}}{スイデン}が\ruby{\underline{見}}{み}\underline{えます}。\\\hline
kawa no DO{\middel}te o a{\middel}garu to SUI{\middel}DEN ga mi{\middel}emasu.\\\hline
(river) (embankment) (go up) (rice paddies) (can see)\\\hline
=\textit{If you go up the river embankment, you can see rice paddies.}\\\hline
\end{sentence}
\end{multicols}
\vhrulefill{2pt}
%%--------------------------------------------------------------------
%%--------------------------------------------------------------------
\begin{multicols}{2}
\begin{glyph}{Hang (掛)}{hang}\label{glyph:hang}
Hang.\\\hline
Note how the last stroke of the bottom ``earth'' component leads into the next line by being drawn with an upward slant.
\end{glyph}

\begin{comps}
\comprtwocone & 扌 + 土 + 土 + ⼘ \\\hline
\comprthreecone & Hand + Earth + Earth + Divination  \\\hline
\end{comps}

\begin{remglyph}
\hooglig{Hanging} \remcom{earth} with your \remcom{hands} is \remcom{divination}.\\\hline
\end{remglyph}

\begin{prons}
\pronsrtwocone & --- & \textbf{か (ka)}\\\hline
\pronsrthreecone & --- & --- \\\hline
\end{prons}

\begin{remprons}
\hooglig{Hanging} clothes is a \comkun{customary} thing to do.\\\hline
\end{remprons}

%{\middel}
% &  ()\\\hline
\begin{somme}
\rye{3}{As the final example below shows, this 漢字 can act similarly to 買 (see entry \ref{glyph:buy}).}
掛かる (intr.) & \textit{to hang (by itself)} & か{\middel}かる \mbox{(ka{\middel}karu)}\\\hline
掛ける (tr.) & \textit{to hang (something)} & か{\middel}ける \mbox{(ka{\middel}keru)}\\\hline
掛金 & hang + gold (money) = \textit{installment} & かけ{\middel}キン \mbox{(kake{\middel}KIN)}\\\hline
\end{somme}

\begin{sentence}
\underline{シャツ}\ruby{を}{お}\underline{あそこ}に\ruby{\underline{掛}}{か}\underline{けて} \ruby{\underline{下}}{くだ}\underline{さい}。\\\hline
SHATSU o asoko ni ka{\middel}kete kuda{\middel}sai.\\\hline
(shirts) (over there) (hang) (please)\\\hline
=\textit{Please hang the shirts over there.}\\\hline
\end{sentence}
\end{multicols}
\vspace*{\fill}
%%--------------------------------------------------------------------
%%--------------------------------------------------------------------
\pagebreak
\begin{multicols}{2}
\begin{glyph}{Ten Thousand (万)}{ten_thousand}\label{glyph:ten_thousand}
Ten thousand.\\\hline
More so than 百 and 千, this character can express the general idea of ``countless amount'' of something.
\end{glyph}

\begin{comps}
\comprtwocone & 丿 + 一 \\\hline
\comprthreecone & NO/Slash + One\\\hline
\end{comps}

\begin{remglyph}
\hooglig{Ten thousand} \remcom{slashes} can be conquered by \remcom{no} \remcom{one}.\\\hline
\end{remglyph}

\begin{prons}
\pronsrtwocone & \textbf{マン (MAN)} & ---\\\hline
\pronsrthreecone & バン (BAN) &  \\\hline
\end{prons}

\begin{remprons}
\hooglig{Ten thousand} \comon{mundane} \ucomon{buns}.\\\hline
\end{remprons}

%{\middel}
% &  ()\\\hline
\begin{somme}
一万 & one + ten thousand = \textit{ten thousand} & イチ{\middel}マン \mbox{(ICHI{\middel}MAN)}\\\hline
二万 & two  + ten thousand = \textit{twenty thousand} & ニ{\middel}マン \mbox{(NI{\middel}MAN)}\\\hline
三万 & three + ten thousand = \textit{thirty thousand} & サン{\middel}マン \mbox{(SAN{\middel}MAN)}\\\hline
四万 & four + ten thousand = \textit{forty thousand} & よん{\middel}マン \mbox{(yon{\middel}MAN)}\\\hline
五万 & five + ten thousand = \textit{fifty thousand} & ゴ{\middel}マン \mbox{(GO{\middel}MAN)}\\\hline
六万 & six + ten thousand = \textit{sixty thousand} & ロク{\middel}マン \mbox{(ROKU{\middel}MAN)}\\\hline
七万 & seven + ten thousand = \textit{seventy thousand} & なな{\middel}マン \mbox{(nana{\middel}MAN)}\\\hline
八万 & eight + ten thousand = \textit{eighty thousand} & ハチ{\middel}マン \mbox{(HACHI{\middel}MAN)}\\\hline
九万 & nine + ten thousand = \textit{ninety thousand} & キュウ{\middel}マン \mbox{(KY\={U}{\middel}MAN)}\\\hline
\end{somme}

\begin{sentence}
\underline{あの}
\ruby{\underline{国}}{くに}の\ruby{\underline{人口}}{ジンコウ}\ruby{は}{わ}\ruby{\underline{八万人}}{ハチマンニン}{\;}\underline{です}。\\\hline
ano kuni no JIN{\middel}K\={O} wa HACHI{\middel}MAN{\middel}NIN desu.\\\hline
(that) (country) (population) (eighty thousand people) (is)\\\hline
=\textit{That country's population is eighty thousand.}\\\hline
\end{sentence}
\end{multicols}
\vhrulefill{2pt}
%%--------------------------------------------------------------------
%%--------------------------------------------------------------------
\begin{multicols}{2}
\begin{glyph}{Rain (雨)}{rain}\label{glyph:rain}
Rain.
\end{glyph}

\begin{comps}
\comprtwocone & 雨 \\\hline
\comprthreecone & Rain \\\hline
\end{comps}

\begin{remglyph}
Part of the 漢字 radicals (\#{173}).\\\hline
\end{remglyph}

\begin{prons}
\pronsrtwocone & --- & \textbf{あめ、あま (ame, ama)}\\\hline
\rye{3}{The 訓読み, あま, only appears in the first position.}
  \pronsrthreecone & ウ (U) & --- \\\hline
\end{prons}

\begin{remprons}
\hooglig{Rain} destroyed the \comon{ooky} \comkun{amalgamation}, and therefore \comkun{amends} were not necessary anymore.\\\hline
\end{remprons}

%{\middel}
% &  ()\\\hline
\begin{somme}
雨 & \textit{rain} & あめ \mbox{(ame)}\\\hline
雨上がり & rain + upper = \textit{after the rain} & あめ　あ{\middel}がり \mbox{(ame a{\middel}gari)}\\\hline
大雨 & large + rain = \textit{heavy rain} & おお{\middel}あめ \mbox{(\={o}{\middel}ame)}\\\hline
\notyet{雨空} & rain + sky = \textit{rainy sky} & あま{\middel}ぞら \mbox{(ama{\middel}zora)}\\\hline
\end{somme}

\begin{irrr}
\rye{3}{The following word can be read with its normal 音読み, but is often heard with the irregular reading pronunciation below.}
    梅雨{\notcov} & plum tree + rain = \textit{the rainy season} & つゆ (tsuyu)\\\hline
\end{irrr}

\begin{sentence}
\ruby{\underline{九月}}{クガツ}に\ruby{は}{わ}\ruby{\underline{日本}}{ニホン}で\ruby{\underline{大雨}}{おおあめ}の\ruby{\underline{日}}{ひ}が\ruby{\underline{多}}{おお}\underline{い}です。\\\hline
KU{\middel}GATSU ni wa NI{\middel}HON de \={o}{\middel}ame no hi ga \={o}{\middel}i desu.\\\hline
(September) (Japan) (heavy rain) (day) (many)\\\hline
=\textit{In September, Japan has many days with heavy rain.}\\\hline
\end{sentence}
\end{multicols}
%%--------------------------------------------------------------------
%%--------------------------------------------------------------------
\pagebreak
\vspace*{\fill}
\begin{multicols}{2}
\begin{glyph}{Electric (電)}{electric}\label{glyph:electric}
Electric; electricity.
\end{glyph}

\begin{comps}
\comprtwocone & 雨 + 日 + 乚 \\\hline
\comprthreecone & Rain + Sun + Second/Latter \\\hline
\end{comps}

\begin{remglyph}
\hooglig{Electricity} is a \remcom{latter} result of \remcom{rain} under the \remcom{sun}.\\\hline
\end{remglyph}

\begin{prons}
\pronsrtwocone & \textbf{デン (DEN)} & ---\\\hline
\pronsrthreecone & --- & --- \\\hline
\end{prons}

\begin{remprons}
\hooglig{Electric} \comon{den}.\\\hline
\end{remprons}

%{\middel}
% &  ()\\\hline
\begin{somme}
    電化 & electric + change = \textit{electrification} & デン{\middel}カ \mbox{(DEN{\middel}KA)}\\\hline
    電力 & electric + strength = \textit{electric power} & ダン{\middel}リョク \mbox{(DEN{\middel}RYOKU)}\\\hline
    電子 & electric + child = \textit{electron} & デン{\middel}シ \mbox{(DEN{\middel}SHI)}\\\hline
    \notyet{電気} & electric + spirit = \textit{electricity} & デン{\middel}キ \mbox{(DEN{\middel}KI)}\\\hline
    \notyet{電話} & electric + speak = \textit{telephone} & デン{\middel}ワ \mbox{(DEN{\middel}WA)}\\\hline
    \notyet{電車} & electric + car = \textit{(electric) train} & デン{\middel}シャ \mbox{(DEN{\middel}SHA)}\\\hline
\end{somme}

\begin{sentence}
\ruby{\underline{田中}}{たなか}\underline{さん}\ruby{は}{わ}\ruby{\underline{電車}}{デンシャ}\ruby{を}{お}\ruby{\underline{下}}{お}\underline{りました}。\\\hline
Ta{\middel}naka-san wa DEN{\middel}SHA o o{\middel}rimashita.\\\hline
(Tanaka-san) (train) (descended)\\\hline
=\textit{Tanaka-san got off the train.}\\\hline
\end{sentence}
\end{multicols}
\vhrulefill{2pt}
%%--------------------------------------------------------------------
%%--------------------------------------------------------------------
\begin{multicols}{2}
\begin{glyph}{Sell (売)}{sell}\label{glyph:sell}
Sell.
\end{glyph}

\begin{comps}
\comprtwocone & 士 + 冖 + 儿\\\hline
    \comprthreecone & Gentleman + Cover/Crown + Man/Person\\\hline
\end{comps}

\begin{remglyph}
By \hooglig{selling} products the \remcom{gentleman} is \remcom{covering} for
    another \remcom{person}.\\\hline
\end{remglyph}

\begin{prons}
\pronsrtwocone & \textbf{バイ (BAI)} & \textbf{う (u)}\\\hline
\pronsrthreecone & --- & --- \\\hline
\end{prons}

\begin{remprons}
    \hooglig{Selling} and \comon{buying} is the \comkun{oomph} behind any business.\\\hline
\end{remprons}

%{\middel}
% &  ()\\\hline
\begin{somme}
\rye{3}{As the final example below shows, this is another kanji that can function similar to 買.}
    売る & \textit{to sell} & う{\middel}る \mbox{(u{\middel}ru)}\\\hline
    売り歩く & sell + walk = \textit{to peddle} & う{\middel}り　ある{\middel}く \mbox{(u{\middel}ri aru{\middel}ku)}\\\hline
    売春 & sell + spring = \textit{prostitution} & バイ{\middel}シュン \mbox{(BAI{\middel}SHUN)}\\\hline
    売り家 & sell + house = \textit{``house for sale''} & う{\middel}り　いえ \mbox{(u{\middel}ri ie)}\\\hline
    \notyet{売り切れる} & sell + cut = \textit{to be sold out} & う{\middel}り　き{\middel}れる \mbox{(u{\middel}ri ki{\middel}reru)}\\\hline
    \notyet{売切れ} & sell + cut = \textit{``sold out''} & う{\middel}り　き{\middel}れ \mbox{(u{\middel}ri ki{\middel}re)}\\\hline
\end{somme}

\begin{sentence}
\ruby{\underline{田口}}{たぐち}\underline{さん}\ruby{は}{わ}\ruby{\underline{日本語}}{ニホンゴ}の\ruby{\underline{本}}{ホン}\ruby{を}{お}\ruby{\underline{売}}{}\underline{って}います。\\\hline
Ta{\middel}guchi-san wa NI{\middel}HON{\middel}GO no HON o u{\middel}tte imasu.\\\hline
(Taguchi-san) (Japanese) (books) (sells)\\\hline
=\textit{Taguchi-san sells Japanese books.}\\\hline
\end{sentence}
\end{multicols}
\vspace*{\fill}
%%--------------------------------------------------------------------
%%--------------------------------------------------------------------
\pagebreak
\begin{multicols}{2}
\begin{glyph}{Individual (者)}{individual}\label{glyph:individual}
Think of individual in the sense of a human being.\\\hline
The characters in a compound that precede this kanji define the person.\\\hline
This character never comes first.
\end{glyph}

\begin{comps}
\comprtwocone & 耂老⺹ + 日 \\\hline
\comprthreecone & Old + Sun \\\hline
\end{comps}

\begin{remglyph}
    That \hooglig{individual} is as \remcom{old} as the \remcom{sun}.\\\hline
\end{remglyph}

\begin{prons}
\pronsrtwocone & \textbf{シャ (SHA)} & ---\\\hline
\pronsrthreecone & --- & もの (mono) \\\hline
\end{prons}

\begin{remprons}
    The \hooglig{individual} \comon{shut} up his \ucomkun{monologue}.\\\hline
\end{remprons}

%{\middel}
% &  ()\\\hline
\begin{somme}
    有力者 & have + strength + individual = \textit{powerful person} & ユウ{\middel}リョク{\middel}シャ \mbox{(Y\={U}{\middel}RYOKU{\middel}SHA)}\\\hline
    \notyet{前者} & before + individual = \textit{the former} & ゼン{\middel}シャ \mbox{(ZEN{\middel}SHA)}\\\hline
    \notyet{後者} & after + individual = \textit{the latter} & コウ{\middel}シャ \mbox{(K\={O}{\middel}SHA)}\\\hline
    \notyet{学者} & study + individual = \textit{scholar} & ガク{\middel}シャ \mbox{(GAKU{\middel}SHA)}\\\hline
\end{somme}

\begin{sentence}
\underline{あの} \ruby{\underline{国}}{くに}で\ruby{\underline{王子}}{オウジ}\ruby{は}{わ}\ruby{\underline{有力者}}{ユウリョクシャ}です。\\\hline
ano kuni de \={O}{\middel}JI wa Y\={U}{\middel}RYOKU{\middel}SHA desu.\\\hline
(that) (country) (prince) (powerful person)\\\hline
=\textit{The prince is a powerful person in that country.}\\\hline
\end{sentence}
\end{multicols}
\vhrulefill{2pt}
%%--------------------------------------------------------------------
%%--------------------------------------------------------------------
\begin{multicols}{2}
\begin{glyph}{Enter (入)}{enter}\label{glyph:enter}
Enter; Putting in.
\end{glyph}

\begin{comps}
\comprtwocone & 入 \\\hline
\comprthreecone & Enter \\\hline
\end{comps}

\begin{remglyph}
Part of the Kanji radicals.\\\hline
\end{remglyph}

\begin{prons}
\pronsrtwocone & \textbf{ニュウ (NY\={U})} & \textbf{はい、い (hai, i)}\\\hline
\pronsrthreecone & --- & --- \\\hline
\end{prons}

\begin{remprons}
\hooglig{Enter} \comon{new} \comkun{interesting} \comkun{heights}.\\\hline
\end{remprons}

%{\middel}
% &  ()\\\hline
\begin{somme}
    入る (intr.) & \textit{to enter} & はい{\middel}る \mbox{(hai{\middel}ru)}\\\hline
    入る (intr.) & \textit{to enter} & い{\middel}る \mbox{(i{\middel}ru)}\\\hline
    入れる (tr.) & \textit{to put (something) into} & い{\middel}れる \mbox{(i{\middel}reru)}\\\hline
    入口 & enter + mouth = \textit{entrance} & いり{\middel}ぐち \mbox{(iri{\middel}guchi)}\\\hline
    入国 & enter + country = \textit{to enter a country} & ニュウ{\middel}コク \mbox{(NY\={U}{\middel}KOKU)}\\\hline
    入金 & enter + gold = \textit{payment} & ニュウ{\middel}キン \mbox{(NY\={U}{\middel}KIN)}\\\hline
    買い入れる & buy + enter = \textit{to stock up on} & か{\middel}い　い{\middel}れる \mbox{(ka{\middel}i i{\middel}reru)}\\\hline
\rye{3}{はい{\middel}る expresses the intransitive meaning of the verb ``enter'' (that is, to enter some type of place).  It is rarely used in compounds.\\\hline
い{\middel}る is most commonly found in compounds, primarily when the other kanji is read with kun-yomi (ニュウ performs this function when the other kanji is read with 音読み).\\\hline
い{\middel}れる is simply the transitive companion of はい{\middel}る.\\\hline
As can be seen in the fourth example, this kanji can behave on rare occasions like 買.}
\end{somme}

\begin{sentence}
\underline{お}\ruby{\underline{金}}{かね}\ruby{を}{お}\ruby{\underline{入}}{い}\underline{れて} \ruby{\underline{下}}{くだ}\underline{さい}。\\\hline
o{\middel}kane o i{\middel}rete kuda{\middel}sai.\\\hline
(money) (put in) (please)\\\hline
=\textit{Please insert the money.}\\\hline
\end{sentence}
\end{multicols}
%%--------------------------------------------------------------------
%%--------------------------------------------------------------------
\pagebreak
\vspace*{\fill}
\begin{multicols}{2}
\begin{glyph}{Snow (雪)}{snow}\label{glyph:snow}
Snow.
\end{glyph}

\begin{comps}
\comprtwocone & 雨 + 彐彑 \\\hline
\comprthreecone & Rain + Pig's head/snout \\\hline
\end{comps}

\begin{remglyph}
\hooglig{Snow} is produced when \remcom{rain} comes out a \remcom{pig's snout}.\\\hline
\end{remglyph}

\begin{prons}
\pronsrtwocone & --- & \textbf{ゆき (yuki)}\\\hline
\pronsrthreecone & セツ (SETSU) & --- \\\hline
\end{prons}

\begin{remprons}
    \hooglig{``Snowy''}, also known as \comkun{``Yuki''}, \comon{sets} up a
    \hooglig{snowman}.\\\hline
\end{remprons}

%{\middel}
% &  ()\\\hline
\begin{somme}
    雪 & \textit{snow} & ゆき \mbox{(yuki)}\\\hline
    雪玉 & snow + jewel = \textit{snowball} & ゆき{\middel}だま \mbox{(yuki{\middel}dama)}\\\hline
    雪女 & snow + woman = \textit{snow fairy} & ゆき{\middel}おんな \mbox{(yuki{\middel}onna)}\\\hline
    雪男 & snow + man = \textit{the Abominable Snowman} & ゆき{\middel}おとｔこ \mbox{(yuki{\middel}otoko)}\\\hline
    大雪 & large + snow = \textit{heavy snow} & おお{\middel}ゆき \mbox{(\={o}{\middel}yuki)}\\\hline
    小雪 & small + snow = \textit{light snow} & こ{\middel}ゆき \mbox{(ko{\middel}yuki)}\\\hline
\end{somme}

\begin{irrr}
    吹雪{\notcov} & blow + snow = \textit{blizzard} & ふぶき (fubuki)\\\hline
    雪崩{\notcov} & snow + crumble = \textit{avalance} & なだれ (nadare)\\\hline
\end{irrr}

\begin{sentence}
\underline{テレビ}の\underline{ニュース}に\underline{よる}と\ruby{\underline{明日}}{あした}\hphantom{h}\underline{から} \ruby{\underline{大雪}}{おおゆき}です。\\\hline
TEREBI no NY\={U}SU ni yoru to ashita kara \={o}yuki desu.\\\hline
(television) (news) (according to) (tomorrow) (from) (heavy snow)\\\hline
=\textit{According to the TV news, heavy snow will be starting tomorrow.}\\\hline
\end{sentence}
\end{multicols}
\vhrulefill{2pt}
%%--------------------------------------------------------------------
%%--------------------------------------------------------------------
\begin{multicols}{2}
\begin{glyph}{Gate (門)}{gate}\label{glyph:gate}
Gate.
\end{glyph}

\begin{comps}
\comprtwocone & 門 \\\hline
\comprthreecone & Gate \\\hline
\end{comps}

\begin{remglyph}
Part of the Kanji radicals (\#{169}).\\\hline
\end{remglyph}

\begin{prons}
\pronsrtwocone & \textbf{モン (MON)} & \textbf{かど (kado)}\\\hline
\pronsrthreecone & --- & --- \\\hline
\end{prons}

\begin{remprons}
    The \hooglig{gates} of \comon{monasteries} are not there for
    \comkun{cuddling}.\\\hline
\end{remprons}

%{\middel}
% &  ()\\\hline
\begin{somme}
    門 & \textit{gate} & モン \mbox{(MON)}\\\hline
    門口 & gate + mouth = \textit{front door} & かど{\middel}ぐち \mbox{(kado{\middel}guchi)}\\\hline
    水門 & water + gate = \textit{sluice gate} & スイ{\middel}モン \mbox{(SUI{\middel}MON)}\\\hline
    校門 & school + gate = \textit{school gate} & コウ{\middel}モン \mbox{(K\={O}{\middel}MON)}\\\hline
\end{somme}

\begin{sentence}
\underline{あそこ}の\ruby{\underline{高}}{たか}\underline{い} \ruby{\underline{門}}{モン}\ruby{は}{わ}\underline{とても} \ruby{\underline{古}}{ふる}\underline{い}ですね。\\\hline
asoko no taka{\middel}i MON wa totemo furu{\middel}i desu ne.\\\hline
(over there) (tall) (gate) (very) (old)\\\hline
=\textit{That tall gate over there is very old, isn't it?}\\\hline
\end{sentence}
\end{multicols}
\vspace*{\fill}
%%--------------------------------------------------------------------
%%--------------------------------------------------------------------
\pagebreak
\vspace*{\fill}
\begin{multicols}{2}
\begin{glyph}{Interval (間)}{interval}\label{glyph:interval}
This kanji indicates an interval in time or space, and thus incorporates the ideas of ``during'', ``between'', and ``among''.\\\hline
When used with the kun-yomi ま, it can also mean a room.
\end{glyph}

\begin{comps}
\comprtwocone & 門 + 日 \\\hline
\comprthreecone & Gate + Sun \\\hline
\end{comps}

\begin{remglyph}
    After an \hooglig{interval} the \remcom{sun} will have passed across the
    \remcom{gate}.\\\hline
\end{remglyph}

\begin{prons}
\pronsrtwocone & \textbf{カン (KAN)} & \textbf{あいだ、ま (aida, ma)}\\\hline
\pronsrthreecone & ケン (KEN) & --- \\\hline
\end{prons}

\begin{remprons}
    During the \hooglig{interval} of the \comon{cunning} \comkun{aida}
    performance, \comkun{muddy} \ucomon{kennels} were set up.\\\hline
\end{remprons}

%{\middel}
% &  ()\\\hline
\begin{somme}
    間 & \textit{interval} & あいだ \mbox{(aida)}\\\hline
    手間 & hand + interval = \textit{labor; trouble} & て{\middel}ま \mbox{(te{\middel}ma)}\\\hline
    中間 & middle + interval = \textit{midway} & チュウ{\middel}カン \mbox{(CH\={U}{\middel}KAN)}\\\hline
    \notyet{時間} & time + interval = \textit{time} & ジ{\middel}カン \mbox{(JI{\middel}KAN)}\\\hline
    \notyet{空間} & empty + interval = \textit{space} & クウ{\middel}カン \mbox{(K\={U}{\middel}KAN)}\\\hline
\end{somme}

\begin{sentence}
\underline{アメリカ}に\underline{いる} \ruby{\underline{間}}{あいだ} \ruby{\underline{英語}}{エイゴ}\ruby{を}{お}\ruby{\underline{話}}{}\underline{します}。\\\hline
AMERIKA ni iru aida EI{\middel}GO o hana{\middel}shimasu.\\\hline
(America) (be) (interval) (English) (speak)\\\hline
=\textit{I'll speak English while I'm in America.}\\\hline
\end{sentence}
\end{multicols}
\vhrulefill{2pt}
%%--------------------------------------------------------------------
%%--------------------------------------------------------------------
\begin{multicols}{2}
\begin{glyph}{Exit (出)}{exit}\label{glyph:exit}
Exit; Leave; Let out.
\end{glyph}

\begin{comps}
\comprtwocone & 屮 + 凵 \\\hline
\comprthreecone & Sprout + Wide opened mouth/Container \\\hline
\end{comps}

\begin{remglyph}
    \hooglig{Exit} the \remcom{sprout} in the \remcom{container}.\\\hline
\end{remglyph}

\begin{prons}
\pronsrtwocone & \textbf{シュツ (SHUTSU)} & \textbf{で、だ (de, da)}\\\hline
\pronsrthreecone & スイ (SUI) & --- \\\hline
\end{prons}

\begin{remprons}
    I made my \hooglig{exit} after I got \comkun{dented} by the \comkun{dull}
    \comon{shoots} in the \ucomon{chop suey}.\\\hline
\end{remprons}

%{\middel}
% &  ()\\\hline
\begin{somme}
    出る (intr.) & \textit{to leave} & で{\middel}る \mbox{(de{\middel}ru)}\\\hline
    出す (tr.) & \textit{to let out} & だ{\middel}す \mbox{(da{\middel}su)}\\\hline
    出かける & \textit{to leave (for somewhere)} & で{\middel}かける \mbox{(de{\middel}kakeru)}\\\hline
    思い出す & think + exit = \textit{to remember} & おも{\middel}い　だ{\middel}す \mbox{(omo{\middel}i da{\middel}su)}\\\hline
    出口 & exit + mouth = \textit{exit} & で{\middel}ぐち \mbox{(de{\middel}guchi)}\\\hline
    出国 & exit + country = \textit{to leave a country} & シュッ{\middel}コク \mbox{(SHUK{\middel}KOKU)}\\\hline
\end{somme}

\begin{sentence}
\underline{すみません}、\ruby{\underline{出口}}{でぐち}\ruby{は}{わ}\underline{どこ}ですか。\\\hline
sumimasen de{\middel}guchi wa doko desu ka.\\\hline
(excuse me) (exit) (where)\\\hline
=\textit{Excuse me, where's the exit?}\\\hline
\end{sentence}
\end{multicols}
\vspace*{\fill}
%%--------------------------------------------------------------------
%%--------------------------------------------------------------------
\pagebreak
\vspace*{\fill}
\begin{multicols}{2}
\begin{glyph}{Life (生)}{life}\label{glyph:life}
One of the most fascinating characters in the language, this kanji deals with life in all its guises.  Along with the ideas of ``birth'', ``production'', and ``growing'', it also incorporates the notions of ``pure'', ``raw'', and ``fresh'', in everything from beer to silk.\\\hline
When used as a suffix, it will often indicate ``student''.
\end{glyph}

\begin{comps}
\comprtwocone & 生 \\\hline
    \comprthreecone & Life \\\hline
\end{comps}

\begin{remglyph}
    Part of the Kanji radicals (\#{100}).\\\hline
\end{remglyph}

\begin{prons}
\pronsrtwocone & \textbf{セイ (SEI)} & \textbf{い、う、なま (i, u, nama)}\\\hline
\rye{3}{なま only appears in the first position.}
\pronsrthreecone & ショウ (SH\={O}) & お、は、き (o, ha, ki) \\\hline
\rye{3}{In terms of variety of pronunciations, kanji do not come any more complicated than this.\\\hline
い and う are verb stems (as are お and は of the less-common 訓読み).}
\end{prons}

\begin{remprons}
    \hooglig{Life} is full \comkun{interesting} \comkun{oomph} when
    \comon{sailing} in \comkun{Namaqualand}.\\\hline
    \hooglig{Life} at the \ucomon{shore} \ucomkun{obligates} me to
    \ucomkun{keep} \ucomkun{halting}.\\\hline
\end{remprons}

%{\middel}
% &  ()\\\hline
\begin{somme}
    生きる (intr.) & \textit{to live} & い{\middel}きる \mbox{(i{\middel}kiru)}\\\hline
    生かす (tr.) & \textit{to give life to} & い{\middel}かす \mbox{(i{\middel}kasu)}\\\hline
    生まれる (intr.) & \textit{to be born} & う{\middel}まれる \mbox{(u{\middel}mareru)}\\\hline
    生む (tr.) & \textit{to give birth to} & う{\middel}む \mbox{(u{\middel}mu)}\\\hline
    生ビール & life + beer = \textit{draft beer} & なま{\middel}ビール \mbox{(nama{\middel}b\={i}ru)}\\\hline
    人生 & person + life = \textit{(human) life} & ジン{\middel}セイ \mbox{(JIN{\middel}SEI)}\\\hline
    \notyet{学生} & study + life = \textit{student} & ガク{\middel}セイ \mbox{(GAKU{\middel}SEI)}\\\hline
    \notyet{先生} & precede + life = \textit{teacher} & セン{\middel}セイ \mbox{(SEN{\middel}SEI)}\\\hline
\end{somme}

\begin{irrr}
    芝生{\notcov} & lawn + life = \textit{lawn} & しば{\middel}ふ (shiba{\middel}fu)\\\hline
\end{irrr}

\begin{sentence}
\ruby{\underline{生}}{なま}\underline{ビール}\ruby{は}{わ}\ruby{\underline{好}}{す}\underline{き}ですか。\\\hline
nama{\middel}b\={i}ru wa su{\middel}ki desu ka.\\\hline
(draft beer) (like)\\\hline
=\textit{Do you like draft beer?}\\\hline
\end{sentence}
\end{multicols}
%%--------------------------------------------------------------------
%%--------------------------------------------------------------------
\begin{multicols}{2}
\begin{glyph}{Cow (牛)}{ushi}\label{glyph:cow}
Cow; Ox; Bull.
\end{glyph}

\begin{comps}
\comprtwocone & 牛 \\\hline
\comprthreecone & Cow \\\hline
\end{comps}

\begin{remglyph}
Part of the Kanji radicals (\#{93}).\\\hline
\end{remglyph}

\begin{prons}
\pronsrtwocone & \textbf{ギュウ (GY\={U})} & \textbf{うし (ushi)}\\\hline
\pronsrthreecone & --- & --- \\\hline
\end{prons}

\begin{remprons}
    \hooglig{Cows} \comkun{shipped} by \comon{gyou}?\\\hline
\end{remprons}

%{\middel}
% &  ()\\\hline
\begin{somme}
    牛 & \textit{cow} & うし \mbox{(ushi)}\\\hline
    子牛 & child + cow = \textit{calf} & こ{\middel}うし \mbox{(ko{\middel}ushi)}\\\hline
    水牛 & water + cow = \textit{water buffalo} & スイ{\middel}ギュウ \mbox{(SUI{\middel}GY\={U})}\\\hline
    牛肉 & cow + meat = \textit{beef} & ギュウ{\middel}ニク \mbox{(GY\={U}{\middel}NIKU)}\\\hline
    \notyet{牛馬} & cow + horse = \textit{cattle and horses} & ギュウ{\middel}バ \mbox{(GY\={U}{\middel}BA)}\\\hline
\end{somme}

\begin{sentence}
\underline{アフリカ}で\ruby{\underline{水牛}}{スイギュウ}と\underline{キリン}\ruby{を}{お}\ruby{\underline{見}}{み}\underline{ました}。\\\hline
AFURIKA de SUI{\middel}GY\={U} to KIRIN o mi{\middel}mashita.\\\hline
(Africa) (water buffalo) (giraffe) (saw)\\\hline
=\textit{(I) saw a water buffalo and giraffe in Africa.}\\\hline
\end{sentence}
\end{multicols}
\vhrulefill{5pt}
\vspace*{\fill}